\documentclass[11pt,a4paper]{article}

\usepackage[margin=2.2cm]{geometry}
\usepackage{amsmath,amssymb,amsthm}
\usepackage{braket}
\usepackage{physics}
\usepackage{booktabs}
\usepackage{enumitem}
\usepackage{xcolor}
\usepackage{array}
\usepackage{tikz}
\usetikzlibrary{quantikz}

\newcommand{\WMC}{\operatorname{WMC}}

\title{Quantum Computing Mock Examination 2\\[0.3em]Complete Solutions}
\author{Summer Semester 2026}
\date{}

\begin{document}
\maketitle
\tableofcontents
\newpage

%% ============================================================
\section{Superposition, Normalization, and Single-Qubit Gates}

\subsection*{(a) Valid quantum states [5 pts]}

A vector $\begin{pmatrix}a\\b\end{pmatrix}$ is a valid quantum state if and only if $a^2+b^2=1$.

\begin{itemize}
  \item[(i)] $(\sqrt{0.6})^2+(\sqrt{0.4})^2=0.6+0.4=1$.\quad
    \textbf{Valid.}\quad $p(0)=0.6$,\; $p(1)=0.4$.
  \item[(ii)] $\bigl(\tfrac{1}{2}\bigr)^2+\bigl(\tfrac{\sqrt{3}}{2}\bigr)^2=\tfrac{1}{4}+\tfrac{3}{4}=1$.\quad
    \textbf{Valid.}\quad $p(0)=\tfrac{1}{4}$,\; $p(1)=\tfrac{3}{4}$.
  \item[(iii)] $1^2+(-1)^2=2\ne1$.\quad \textbf{Not valid.}
  \item[(iv)] $(-\sqrt{0.75})^2+(\sqrt{0.25})^2=0.75+0.25=1$.\quad
    \textbf{Valid.}\quad $p(0)=0.75$,\; $p(1)=0.25$.
\end{itemize}

\noindent\textbf{Key observation:} States (i) and (iv) have different amplitude vectors but yield
the same probability distribution. Signs of amplitudes do not affect measurement probabilities;
they matter only for \emph{interference} effects.

\subsection*{(b) Gate sequence $H\to Z\to H$ and the identity $HZH=X$ [5 pts]}

Recall $H=\frac{1}{\sqrt{2}}\begin{pmatrix}1&1\\1&-1\end{pmatrix}$
and $Z=\begin{pmatrix}1&0\\0&-1\end{pmatrix}$.

\medskip
Starting from $\ket{0}$:

\begin{enumerate}[label=\arabic*.]
  \item $H\ket{0}=\dfrac{1}{\sqrt{2}}\begin{pmatrix}1\\1\end{pmatrix}=:\ket{+}$.
  \item $Z\ket{+}=\dfrac{1}{\sqrt{2}}\begin{pmatrix}1&0\\0&-1\end{pmatrix}\begin{pmatrix}1\\1\end{pmatrix}
    =\dfrac{1}{\sqrt{2}}\begin{pmatrix}1\\-1\end{pmatrix}=:\ket{-}$.
  \item $H\ket{-}=\dfrac{1}{\sqrt{2}}\begin{pmatrix}1&1\\1&-1\end{pmatrix}\dfrac{1}{\sqrt{2}}\begin{pmatrix}1\\-1\end{pmatrix}
    =\dfrac{1}{2}\begin{pmatrix}0\\2\end{pmatrix}=\begin{pmatrix}0\\1\end{pmatrix}=\ket{1}$.
\end{enumerate}

Direct matrix computation confirms the identity:
\[
HZH=\frac{1}{2}\begin{pmatrix}1&1\\1&-1\end{pmatrix}\begin{pmatrix}1&0\\0&-1\end{pmatrix}\begin{pmatrix}1&1\\1&-1\end{pmatrix}
=\frac{1}{2}\begin{pmatrix}1&-1\\1&1\end{pmatrix}\begin{pmatrix}1&1\\1&-1\end{pmatrix}
=\begin{pmatrix}0&1\\1&0\end{pmatrix}=X.
\]

\noindent\textbf{Conclusion:} The composite gate $HZH=X$ (Pauli $X$ / NOT gate). This is a fundamental
circuit identity: conjugating $Z$ by Hadamard layers yields $X$, and dually $HXH=Z$.

\subsection*{(c) Measurement with and without Hadamard [5 pts]}

State: $\ket{\psi}=\frac{1}{2}\ket{0}+\frac{\sqrt{3}}{2}\ket{1}$.

\textbf{(i)} Probability of measuring $\ket{0}$: $p(0)=\bigl(\tfrac{1}{2}\bigr)^2=\dfrac{1}{4}$.
Post-measurement state: collapses to $\ket{0}$ with certainty.

\textbf{(ii)} Apply $H$ first:
\begin{align*}
H\ket{\psi} &= \tfrac{1}{2}H\ket{0}+\tfrac{\sqrt{3}}{2}H\ket{1}
= \tfrac{1}{2}\cdot\tfrac{1}{\sqrt{2}}\begin{pmatrix}1\\1\end{pmatrix}
  +\tfrac{\sqrt{3}}{2}\cdot\tfrac{1}{\sqrt{2}}\begin{pmatrix}1\\-1\end{pmatrix}
= \tfrac{1}{2\sqrt{2}}\begin{pmatrix}1+\sqrt{3}\\1-\sqrt{3}\end{pmatrix}.
\end{align*}
\[
p(0)=\left(\tfrac{1+\sqrt{3}}{2\sqrt{2}}\right)^2=\tfrac{(1+\sqrt{3})^2}{8}
=\tfrac{4+2\sqrt{3}}{8}=\tfrac{2+\sqrt{3}}{4}\approx0.933.
\]
\[
p(1)=\left(\tfrac{1-\sqrt{3}}{2\sqrt{2}}\right)^2=\tfrac{4-2\sqrt{3}}{8}=\tfrac{2-\sqrt{3}}{4}\approx0.067.
\]
Check: $\tfrac{2+\sqrt{3}}{4}+\tfrac{2-\sqrt{3}}{4}=1$.\;$\checkmark$

\subsection*{(d) Orthogonality and self-inverse property of $X$ and $Z$ [5 pts]}

\textbf{Orthogonality ($U^TU=I$):}

For $X$: $X^T=X$ (symmetric), so $X^TX=X^2=\begin{pmatrix}0&1\\1&0\end{pmatrix}^2=\begin{pmatrix}1&0\\0&1\end{pmatrix}=I$.\;$\checkmark$

For $Z$: $Z^T=Z$ (diagonal $\Rightarrow$ symmetric), so $Z^TZ=Z^2=\begin{pmatrix}1&0\\0&1\end{pmatrix}=I$.\;$\checkmark$

\textbf{Self-inverse ($U^2=I$):} Follows immediately from the same computations above.

\textbf{Physical significance:}
\begin{itemize}
  \item \emph{Orthogonality} guarantees norm preservation: $\|U\vec{v}\|^2=\vec{v}^TU^TU\vec{v}=\|\vec{v}\|^2$,
    ensuring a normalized quantum state remains normalized after the gate is applied.
    This is a fundamental requirement: total probability must remain 1.
  \item \emph{Self-inverse} ($U^2=I$) means the gate is its own inverse, enabling straightforward
    \emph{uncomputation} in reversible circuits. Applying the same gate twice cancels it,
    which is exploited in ancilla cleanup strategies (Bennett's trick).
\end{itemize}

\newpage
%% ============================================================
\section{Kronecker Products and Parallel Composition}

Given:
\[
\mathbf{P}_A=\begin{pmatrix}\tfrac34&\tfrac14\\[4pt]\tfrac12&\tfrac12\end{pmatrix},\qquad
\mathbf{P}_B=\begin{pmatrix}\tfrac23&\tfrac13\\[4pt]\tfrac14&\tfrac34\end{pmatrix}.
\]

\subsection*{(a) Joint transition matrix $\mathbf{P}_{AB}=\mathbf{P}_A\otimes\mathbf{P}_B$ [5 pts]}

State ordering: $00,01,10,11$. Entry $(i_A,i_B)\to(j_A,j_B)$: $P_A(i_A,j_A)\cdot P_B(i_B,j_B)$.

\medskip
\begin{center}
\begin{tabular}{c|cccc}
row $\backslash$ col & $00$ & $01$ & $10$ & $11$ \\
\hline
$00$ & $\tfrac34\cdot\tfrac23=\tfrac12$ & $\tfrac34\cdot\tfrac13=\tfrac14$ & $\tfrac14\cdot\tfrac23=\tfrac16$ & $\tfrac14\cdot\tfrac13=\tfrac{1}{12}$ \\[4pt]
$01$ & $\tfrac34\cdot\tfrac14=\tfrac{3}{16}$ & $\tfrac34\cdot\tfrac34=\tfrac{9}{16}$ & $\tfrac14\cdot\tfrac14=\tfrac{1}{16}$ & $\tfrac14\cdot\tfrac34=\tfrac{3}{16}$ \\[4pt]
$10$ & $\tfrac12\cdot\tfrac23=\tfrac13$ & $\tfrac12\cdot\tfrac13=\tfrac16$ & $\tfrac12\cdot\tfrac23=\tfrac13$ & $\tfrac12\cdot\tfrac13=\tfrac16$ \\[4pt]
$11$ & $\tfrac12\cdot\tfrac14=\tfrac18$ & $\tfrac12\cdot\tfrac34=\tfrac38$ & $\tfrac12\cdot\tfrac14=\tfrac18$ & $\tfrac12\cdot\tfrac34=\tfrac38$ \\
\end{tabular}
\end{center}

\[
\mathbf{P}_{AB}=
\begin{pmatrix}
\tfrac{1}{2}  & \tfrac{1}{4}  & \tfrac{1}{6}  & \tfrac{1}{12} \\[5pt]
\tfrac{3}{16} & \tfrac{9}{16} & \tfrac{1}{16} & \tfrac{3}{16} \\[5pt]
\tfrac{1}{3}  & \tfrac{1}{6}  & \tfrac{1}{3}  & \tfrac{1}{6}  \\[5pt]
\tfrac{1}{8}  & \tfrac{3}{8}  & \tfrac{1}{8}  & \tfrac{3}{8}
\end{pmatrix}
\]

Row sums: $\tfrac{6+3+2+1}{12}=1$;\; $\tfrac{3+9+1+3}{16}=1$;\; $\tfrac{2+1+2+1}{6}=1$;\; $\tfrac{1+3+1+3}{8}=1$.\;$\checkmark$

\subsection*{(b) Distribution after one step [5 pts]}

Starting in state $00$: $\vec{p}(0)=(1,0)\otimes(1,0)=(1,0,0,0)$.

\textbf{Method 1 (direct):} $\vec{p}(1)=\vec{p}(0)\mathbf{P}_{AB}=$ row $00$ of $\mathbf{P}_{AB}$:
\[
\vec{p}(1)=\Bigl(\tfrac{1}{2},\;\tfrac{1}{4},\;\tfrac{1}{6},\;\tfrac{1}{12}\Bigr).
\]

\textbf{Method 2 (separate chains):}
\[
\vec{q}_A(1)=(1,0)\mathbf{P}_A=\Bigl(\tfrac{3}{4},\tfrac{1}{4}\Bigr),\qquad
\vec{q}_B(1)=(1,0)\mathbf{P}_B=\Bigl(\tfrac{2}{3},\tfrac{1}{3}\Bigr).
\]
\[
\vec{q}_A(1)\otimes\vec{q}_B(1)=\Bigl(\tfrac{3}{4}\cdot\tfrac{2}{3},\;\tfrac{3}{4}\cdot\tfrac{1}{3},\;\tfrac{1}{4}\cdot\tfrac{2}{3},\;\tfrac{1}{4}\cdot\tfrac{1}{3}\Bigr)=\Bigl(\tfrac{1}{2},\;\tfrac{1}{4},\;\tfrac{1}{6},\;\tfrac{1}{12}\Bigr).\;\checkmark
\]

\subsection*{(c) Stationary distribution of the joint chain [5 pts]}

\textbf{Chain $A$} (solve $\vec{\pi}_A\mathbf{P}_A=\vec{\pi}_A$, col $0$):
\[
\pi_{A,0}=\tfrac{3}{4}\pi_{A,0}+\tfrac{1}{2}\pi_{A,1}
\;\Rightarrow\;\tfrac{1}{4}\pi_{A,0}=\tfrac{1}{2}\pi_{A,1}
\;\Rightarrow\;\pi_{A,0}=2\pi_{A,1}.
\]
Normalization: $3\pi_{A,1}=1$, so $\pi_{A,1}=\tfrac{1}{3}$, $\pi_{A,0}=\tfrac{2}{3}$.
\[\vec{\pi}_A=\Bigl(\tfrac{2}{3},\;\tfrac{1}{3}\Bigr).\]

\textbf{Chain $B$} (col $0$):
\[
\pi_{B,0}=\tfrac{2}{3}\pi_{B,0}+\tfrac{1}{4}\pi_{B,1}
\;\Rightarrow\;\tfrac{1}{3}\pi_{B,0}=\tfrac{1}{4}\pi_{B,1}
\;\Rightarrow\;\pi_{B,1}=\tfrac{4}{3}\pi_{B,0}.
\]
Normalization: $\pi_{B,0}\bigl(1+\tfrac{4}{3}\bigr)=\pi_{B,0}\cdot\tfrac{7}{3}=1$,
so $\pi_{B,0}=\tfrac{3}{7}$, $\pi_{B,1}=\tfrac{4}{7}$.
\[\vec{\pi}_B=\Bigl(\tfrac{3}{7},\;\tfrac{4}{7}\Bigr).\]

\textbf{Joint:}
\[
\vec{\pi}_{AB}=\vec{\pi}_A\otimes\vec{\pi}_B
=\Bigl(\tfrac{2}{3}\cdot\tfrac{3}{7},\;\tfrac{2}{3}\cdot\tfrac{4}{7},\;\tfrac{1}{3}\cdot\tfrac{3}{7},\;\tfrac{1}{3}\cdot\tfrac{4}{7}\Bigr)
=\Bigl(\tfrac{2}{7},\;\tfrac{8}{21},\;\tfrac{1}{7},\;\tfrac{4}{21}\Bigr).
\]

\textbf{Proof via mixed-product property:}
\[
\vec{\pi}_{AB}\,\mathbf{P}_{AB}
=(\vec{\pi}_A\otimes\vec{\pi}_B)(\mathbf{P}_A\otimes\mathbf{P}_B)
\overset{\text{MPP}}{=}(\vec{\pi}_A\mathbf{P}_A)\otimes(\vec{\pi}_B\mathbf{P}_B)
=\vec{\pi}_A\otimes\vec{\pi}_B=\vec{\pi}_{AB}.\quad\square
\]

\subsection*{(d) Kronecker product in a quantum circuit: $M'=U\otimes I_2$ [5 pts]}

Given $U=\begin{pmatrix}\sqrt{0.75}&\sqrt{0.25}\\\sqrt{0.25}&-\sqrt{0.75}\end{pmatrix}$,
the $4\times4$ matrix for $U$ acting on the top qubit (identity on bottom) is:
\[
M'=U\otimes I_2=
\begin{pmatrix}
\sqrt{0.75}&0&\sqrt{0.25}&0\\
0&\sqrt{0.75}&0&\sqrt{0.25}\\
\sqrt{0.25}&0&-\sqrt{0.75}&0\\
0&\sqrt{0.25}&0&-\sqrt{0.75}
\end{pmatrix}.
\]

For input $\ket{11}=(0,0,0,1)^T$ (index 3), read off column 3 of $M'$:
\[
M'\ket{11}=\begin{pmatrix}0\\\sqrt{0.25}\\0\\-\sqrt{0.75}\end{pmatrix}
=\sqrt{0.25}\,\ket{01}-\sqrt{0.75}\,\ket{11}.
\]

Measurement probabilities: $p(01)=(\sqrt{0.25})^2=0.25$;\quad $p(11)=(-\sqrt{0.75})^2=0.75$.
Sum: $1$.\;$\checkmark$

\noindent\textit{Interpretation:} $U$ acts on the top qubit via $U\ket{1}=\sqrt{0.25}\ket{0}-\sqrt{0.75}\ket{1}$,
while the bottom qubit $\ket{1}$ is unchanged --- exactly what $U\otimes I$ computes.

\newpage
%% ============================================================
\section{Entanglement and Bell States}

\subsection*{(a) Preparing $\ket{\Phi^+}$ from $\ket{00}$ [5 pts]}

Circuit: apply $H$ to qubit 1, then $\mathrm{CNOT}_{1\to2}$.

\[
\begin{quantikz}
\lstick{$\ket{0}$} & \gate{H} & \ctrl{1} & \rstick[2]{$\ket{\Phi^+}$}\qw \\
\lstick{$\ket{0}$} & \qw      & \targ{}  & \qw
\end{quantikz}
\]

\textbf{Step 1} ($H\otimes I$):
\[
(H\otimes I)\ket{00}=H\ket{0}\otimes\ket{0}=\tfrac{1}{\sqrt{2}}(\ket{0}+\ket{1})\otimes\ket{0}
=\tfrac{1}{\sqrt{2}}(\ket{00}+\ket{10}).
\]

\textbf{Step 2} ($\mathrm{CNOT}_{1\to2}$, mapping $\ket{x,y}\to\ket{x,x\oplus y}$):
\[
\tfrac{1}{\sqrt{2}}(\ket{00}+\ket{10})\;\longrightarrow\;
\tfrac{1}{\sqrt{2}}(\ket{00}+\ket{11})=\ket{\Phi^+}.\quad\checkmark
\]

\subsection*{(b) $\ket{\Phi^+}$ is entangled [5 pts]}

Suppose $\ket{\Phi^+}=(\alpha_1\ket{0}+\beta_1\ket{1})\otimes(\alpha_2\ket{0}+\beta_2\ket{1})$.
Expanding:
\[
=\alpha_1\alpha_2\ket{00}+\alpha_1\beta_2\ket{01}+\beta_1\alpha_2\ket{10}+\beta_1\beta_2\ket{11}.
\]
Matching with $\ket{\Phi^+}=\tfrac{1}{\sqrt{2}}\ket{00}+0\cdot\ket{01}+0\cdot\ket{10}+\tfrac{1}{\sqrt{2}}\ket{11}$:
\begin{align}
\alpha_1\alpha_2 &= \tfrac{1}{\sqrt{2}} \ne 0 \tag{1}\\
\alpha_1\beta_2  &= 0 \tag{2}\\
\beta_1\alpha_2  &= 0 \tag{3}\\
\beta_1\beta_2   &= \tfrac{1}{\sqrt{2}} \ne 0 \tag{4}
\end{align}
From (1): $\alpha_1\ne0$ and $\alpha_2\ne0$.\;
From (2) with $\alpha_1\ne0$: $\beta_2=0$.\;
But (4): $\beta_1\cdot0=\tfrac{1}{\sqrt{2}}\ne0$.\quad\textbf{Contradiction.}\quad$\square$

\subsection*{(c) Partial measurement of $\ket{\Phi^+}$ [5 pts]}

Write $\ket{\Phi^+}=\tfrac{1}{\sqrt{2}}\ket{0}_1\otimes\ket{0}_2+\tfrac{1}{\sqrt{2}}\ket{1}_1\otimes\ket{1}_2$.

\begin{itemize}
  \item \textbf{Outcome 0:}\quad $p(0)=\bigl(\tfrac{1}{\sqrt{2}}\bigr)^2=\tfrac{1}{2}$.
    Post-measurement full state (renormalized): $\ket{00}$.
    Second qubit: $\ket{0}$.
  \item \textbf{Outcome 1:}\quad $p(1)=\tfrac{1}{2}$.
    Post-measurement full state: $\ket{11}$.
    Second qubit: $\ket{1}$.
\end{itemize}

\noindent\textbf{Physical significance:} The two qubits are perfectly correlated.
Measuring qubit 1 as $0$ instantly forces qubit 2 into $\ket{0}$, and vice versa, regardless
of the spatial separation between the qubits. This non-classical correlation without communication is
the defining feature of quantum entanglement (the Einstein-Podolsky-Rosen phenomenon).

\subsection*{(d) All Bell states from $\ket{\Phi^+}$ via local operations [5 pts]}

Apply $U\otimes I$ to $\ket{\Phi^+}=\tfrac{1}{\sqrt{2}}(\ket{00}+\ket{11})$:
$(U\otimes I)\ket{\Phi^+}=\tfrac{1}{\sqrt{2}}(U\ket{0}\otimes\ket{0}+U\ket{1}\otimes\ket{1})$.

\medskip
\begin{center}
\begin{tabular}{l|l|l|l}
Operation & $U\ket{0}$ & $U\ket{1}$ & Result \\\hline
$I\otimes I$ & $\ket{0}$ & $\ket{1}$ & $\tfrac{1}{\sqrt{2}}(\ket{00}+\ket{11})=\ket{\Phi^+}$ \\[4pt]
$Z\otimes I$ & $\ket{0}$ & $-\ket{1}$ & $\tfrac{1}{\sqrt{2}}(\ket{00}-\ket{11})=\ket{\Phi^-}$ \\[4pt]
$X\otimes I$ & $\ket{1}$ & $\ket{0}$ & $\tfrac{1}{\sqrt{2}}(\ket{10}+\ket{01})=\ket{\Psi^+}$ \\[4pt]
$XZ\otimes I$ & $X\ket{0}=\ket{1}$ & $X(-\ket{1})=-\ket{0}$ & $\tfrac{1}{\sqrt{2}}(\ket{10}-\ket{01})=-\ket{\Psi^-}\equiv\ket{\Psi^-}$
\end{tabular}
\end{center}

\noindent All four Bell states are reached.\;$\checkmark$\;
This is the key principle behind \textbf{superdense coding}: two classical bits of information
can be encoded into a single qubit by choosing one of the four local operations on one half of an
entangled Bell pair, then retrieved by a Bell-basis measurement.

\newpage
%% ============================================================
\section{The Deutsch-Jozsa Algorithm}

\subsection*{(a) $(H^{\otimes n})^2=I_{2^n}$ [4 pts]}

First verify $H^2=I_2$:
\[
H^2=\tfrac{1}{2}\begin{pmatrix}1&1\\1&-1\end{pmatrix}\begin{pmatrix}1&1\\1&-1\end{pmatrix}
=\tfrac{1}{2}\begin{pmatrix}2&0\\0&2\end{pmatrix}=I_2.
\]
Then by the mixed-product property $(A\otimes B)(C\otimes D)=(AC)\otimes(BD)$:
\[
(H^{\otimes n})^2
=(H\otimes\cdots\otimes H)(H\otimes\cdots\otimes H)
=(H^2)\otimes\cdots\otimes(H^2)
=I_2\otimes\cdots\otimes I_2=I_{2^n}.\quad\square
\]
Consequence: $H^{\otimes n}$ is self-inverse. The final Hadamard in the DJ circuit exactly
undoes the opening superposition step, enabling a \emph{deterministic} zero-error outcome.

\subsection*{(b) DJ for $n=2$, constant function $f\equiv1$ [5 pts]}

\textbf{Initial:} $\ket{00}\ket{1}$.

\textbf{After $H^{\otimes3}$:}
\[
\tfrac{1}{2}(\ket{00}+\ket{01}+\ket{10}+\ket{11})\otimes\ket{-}
=\tfrac{1}{2}\sum_{x\in\{0,1\}^2}\ket{x}\otimes\ket{-}.
\]

\textbf{After oracle} ($f\equiv1$, so $(-1)^{f(x)}=-1$ for all $x$):
\[
-\tfrac{1}{2}\sum_x\ket{x}\otimes\ket{-}\;\equiv\;\tfrac{1}{2}\sum_x\ket{x}\otimes\ket{-}
\quad(\text{global phase }-1\text{ is unobservable}).
\]
The query register remains $\ket{+}\ket{+}$.

\textbf{After $H^{\otimes2}$ on query register} ($H\ket{+}=\ket{0}$):
\[
\ket{+}\ket{+}\;\xrightarrow{H\otimes H}\;\ket{0}\ket{0}=\ket{00}.
\]

\textbf{Measurement:} $\ket{00}$ with probability $1$.\;
Since result $=\ket{00}$, conclude $f$ is \textbf{constant}.\;$\checkmark$

\subsection*{(c) DJ for $n=2$, balanced function $f(x_1,x_2)=x_1\oplus x_2$ [5 pts]}

Function values: $f(00)=0,\;f(01)=1,\;f(10)=1,\;f(11)=0$. Balanced.\;$\checkmark$

\textbf{After $H^{\otimes3}$:} same as (b), $\tfrac{1}{2}\sum_x\ket{x}\otimes\ket{-}$.

\textbf{After oracle} (phases $(-1)^{f(x)}$):
\[
\tfrac{1}{2}\bigl[\ket{00}-\ket{01}-\ket{10}+\ket{11}\bigr]\otimes\ket{-}.
\]
Factor the query register:
\[
\ket{00}-\ket{01}-\ket{10}+\ket{11}
=(\ket{0}-\ket{1})\otimes(\ket{0}-\ket{1})
=\bigl(\sqrt{2}\,\ket{-}\bigr)\otimes\bigl(\sqrt{2}\,\ket{-}\bigr)
=2\,\ket{-}\otimes\ket{-}.
\]
So the (normalized) query register is $\ket{-}\otimes\ket{-}$.

\textbf{After $H^{\otimes2}$} ($H\ket{-}=\ket{1}$):
\[
\ket{-}\otimes\ket{-}\;\xrightarrow{H\otimes H}\;\ket{1}\ket{1}=\ket{11}.
\]

\textbf{Measurement:} $\ket{11}$ with probability $1$.
Since result $\ne\ket{00}$, conclude $f$ is \textbf{balanced}.\;$\checkmark$

\subsection*{(d) Classical lower bound: $2^{n-1}+1$ queries required [6 pts]}

\textbf{Claim:} $2^{n-1}$ queries are insufficient for a deterministic algorithm.

\textbf{Proof (adversary argument):}
Suppose the algorithm has made $k\le2^{n-1}$ queries and observed all output values equal to $0$
(the all-$1$ case is symmetric).

\begin{itemize}
  \item \emph{Consistent constant function:} $f\equiv0$. All $k$ observed outputs match.\;$\checkmark$
  \item \emph{Consistent balanced function:} Set $f(x)=0$ on the $k$ queried inputs
    and $f(x)=1$ on exactly $2^{n-1}$ of the $2^n-k\ge2^{n-1}$ unqueried inputs
    (and $0$ on the rest). Since $2^n-k\ge2^{n-1}$, there are enough unqueried inputs
    to accommodate $2^{n-1}$ ones.\;$\checkmark$
\end{itemize}

Since both a constant and a balanced function are consistent with the observed outputs,
no deterministic algorithm can distinguish them. Hence $k\le2^{n-1}$ queries are insufficient.

\textbf{Sufficiency of $2^{n-1}+1$ queries:}
After $2^{n-1}+1$ queries all returning $0$, the remaining unqueried inputs number
$2^n-(2^{n-1}+1)=2^{n-1}-1$.
A balanced function needs exactly $2^{n-1}$ ones in total; but only $2^{n-1}-1<2^{n-1}$ unqueried
inputs remain to hold them. Contradiction. So the function must be $f\equiv0$ (constant).\;$\square$

\textbf{Conclusion:} Classical deterministic query complexity $=2^{n-1}+1$ vs.\ quantum complexity $=1$
--- an \textbf{exponential} separation.

\newpage
%% ============================================================
\section{Algebraic Normal Form (ANF) and Circuit Optimization}

\subsection*{(a) Elementary operations in ANF [5 pts]}

\begin{enumerate}[label=(\roman*)]
  \item $\neg A = 1\oplus A$.\quad Verification: $A=0\Rightarrow1\oplus0=1=\neg0$;\;
    $A=1\Rightarrow1\oplus1=0=\neg1$.\;$\checkmark$

  \item $A\lor B = A\oplus B\oplus AB$.\quad Verification:
    \begin{center}
    \begin{tabular}{cc|c|c}
    $A$&$B$&$A\lor B$&$A\oplus B\oplus AB$\\\hline
    0&0&0&$0\oplus0\oplus0=0$\;$\checkmark$\\
    0&1&1&$0\oplus1\oplus0=1$\;$\checkmark$\\
    1&0&1&$1\oplus0\oplus0=1$\;$\checkmark$\\
    1&1&1&$1\oplus1\oplus1=1$\;$\checkmark$
    \end{tabular}
    \end{center}

  \item $\neg(AB) = 1\oplus AB$.\quad Follows from (i) with $A\cdot B$ substituted.
    $A=B=1$: $1\oplus1=0=\neg1$.\;$A=1,B=0$: $1\oplus0=1=\neg0$.\;$\checkmark$
\end{enumerate}

\subsection*{(b) ANF of $f(A,B,C)=(A\lor B)\land C$ [5 pts]}

\textbf{Step 1:} From (a)(ii), $A\lor B=A\oplus B\oplus AB$.

\textbf{Step 2:} Multiply by $C$ using distributivity $(X\oplus Y)\cdot Z=XZ\oplus YZ$:
\[
f=(A\oplus B\oplus AB)\cdot C=AC\oplus BC\oplus ABC.
\]

\[
\boxed{f(A,B,C)=AC\oplus BC\oplus ABC.}
\]

\textbf{Verification (all 8 inputs):}
\begin{center}
\begin{tabular}{ccc|c|c}
$A$&$B$&$C$&$(A\lor B)\land C$&$AC\oplus BC\oplus ABC$\\\hline
0&0&0&0&$0\oplus0\oplus0=0$\;$\checkmark$\\
0&0&1&0&$0\oplus0\oplus0=0$\;$\checkmark$\\
0&1&0&0&$0\oplus0\oplus0=0$\;$\checkmark$\\
0&1&1&1&$0\oplus1\oplus0=1$\;$\checkmark$\\
1&0&0&0&$0\oplus0\oplus0=0$\;$\checkmark$\\
1&0&1&1&$1\oplus0\oplus0=1$\;$\checkmark$\\
1&1&0&0&$0\oplus0\oplus0=0$\;$\checkmark$\\
1&1&1&1&$1\oplus1\oplus1=1$\;$\checkmark$
\end{tabular}
\end{center}

\subsection*{(c) Reversible circuit from ANF [5 pts]}

ANF: $f=AC\oplus BC\oplus ABC$. Goal: $y\leftarrow y\oplus f(A,B,C)$ with $y$ initially $0$.

Each degree-$d$ monomial in the ANF contributes one gate:
degree 2 $\to$ Toffoli (2 controls, 1 target);
degree 3 $\to$ requires an ancilla scratch qubit $s$ (initially $0$).

\medskip
\begin{enumerate}[label=\arabic*.]
  \item $\operatorname{Toffoli}(A,C;\,y)$:\quad $y\leftarrow y\oplus AC$.
  \item $\operatorname{Toffoli}(B,C;\,y)$:\quad $y\leftarrow y\oplus BC$.
  \item $\operatorname{Toffoli}(A,B;\,s)$:\quad $s\leftarrow AB$ \hfill (scratch ancilla)
  \item $\operatorname{Toffoli}(s,C;\,y)$:\quad $y\leftarrow y\oplus ABC$.
  \item $\operatorname{Toffoli}(A,B;\,s)$:\quad $s\leftarrow0$ \hfill (uncompute $s$)
\end{enumerate}

\textbf{Total:} 5 Toffoli gates, 0 CNOT gates, 1 ancilla qubit $s$.
After step 5, $s=0$ is restored. All original inputs $A,B,C$ are preserved on their wires,
confirming reversibility.

\subsection*{(d) Comparison with naive minterm circuit [5 pts]}

\textbf{(i) Minterms of $f$} (where $(A\lor B)\land C=1$, i.e.\ $C=1$ and at least one of $A,B$ is $1$):
\[
(A,B,C)\in\{(0,1,1),\;(1,0,1),\;(1,1,1)\}.\quad\textbf{3 minterms.}
\]

\textbf{(ii) Naive gate count:} Each minterm needs one 3-controlled NOT on $y$,
implemented with NOT gates on $0$-valued control bits:
\begin{itemize}
  \item $(0,1,1)$: $X_A$, Toffoli$(A,B,C;\,y)$, $X_A$.\quad$=1$ Toffoli$+2$ NOTs.
  \item $(1,0,1)$: $X_B$, Toffoli$(A,B,C;\,y)$, $X_B$.\quad$=1$ Toffoli$+2$ NOTs.
  \item $(1,1,1)$: Toffoli$(A,B,C;\,y)$ with ancilla decomposition.\quad$\approx 2$ Toffoli$+1$ ancilla.
\end{itemize}
Approximate total: $\approx4$ Toffoli$+4$ NOT gates.

\textbf{(iii) Scaling:}

With $n$ input bits, a function can have up to $2^{n-1}$ minterms.
Each degree-$n$ minterm requires $O(n)$ elementary gates to decompose.
Naive approach: up to $O(n\cdot 2^n)$ gates.

ANF approach: the number of ANF terms is at most $2^n$, but structured Boolean functions
(SAT oracles, physical Hamiltonians) often have compact ANFs with $O(\text{poly}(n))$ terms.
Each degree-$d$ ANF term requires $O(d)$ Toffoli gates.
For such functions the ANF circuit is polynomial-size, vs.\ exponential for the naive minterm circuit.

Furthermore, the ANF does not require additional NOT gates for negated controls.
In the NISQ context, this compact encoding has reduced ancilla qubit requirements from
$O(n)$ to a single ancilla qubit in practical Grover oracle benchmarks (see Lecture 10).

\newpage
%% ============================================================
\section{Quantum Complexity and Counting Arguments}

\subsection*{(a) Counting argument [6 pts]}

\textbf{Number of Boolean functions:}
$f:\{0,1\}^n\to\{0,1\}$ is determined by $2^n$ independent binary values,
giving $2^{2^n}$ possible functions (doubly exponential in $n$).

\textbf{Number of circuits with $\le C$ gates:}
Fix a gate set $\mathcal{G}$ where each gate acts on at most $k$ qubits.
The number of ways to choose one gate is $|\mathcal{G}|$ (type) times $\binom{n}{k}\le n^k$ (qubit positions).
With $C$ sequential gates:
\[
\#\{\text{circuits with}\le C\text{ gates}\}\;\le\;\bigl(|\mathcal{G}|\cdot n^k\bigr)^C.
\]

\textbf{Gap:}
For a circuit to compute a specific function exactly, distinct functions require distinct circuits.
To cover even a fraction $1/100$ of all $2^{2^n}$ functions:
\[
(|\mathcal{G}|\cdot n^k)^C\ge\tfrac{1}{100}\cdot2^{2^n}
\quad\Rightarrow\quad
C\cdot\log(|\mathcal{G}|\cdot n^k)\ge2^n-O(1)
\quad\Rightarrow\quad
C=\Omega\!\left(\frac{2^n}{n}\right).
\]

\textbf{Conclusion:} The overwhelming majority of Boolean functions require at least $\Omega(2^n/n)$
quantum gates to compute --- \emph{exponentially many}.
Quantum computers, like classical ones, can only compute an exponentially small fraction of all
functions efficiently. Quantum speedups arise only for \emph{specific structured problems}
(factoring, search, quantum simulation).

\subsection*{(b) EQP vs.\ BQP [4 pts]}

\textbf{EQP (Exact Quantum Polynomial time):}
Problems solvable in polynomial time by a quantum computer with \emph{zero} error probability;
the algorithm must always give the correct answer.

\textbf{BQP (Bounded-Error Quantum Polynomial time):}
Problems solvable in polynomial time by a quantum computer with error probability $\le\tfrac{1}{3}$.
The error bound can be amplified to any $\varepsilon>0$ by running $O(\log(1/\varepsilon))$ repetitions.

\begin{enumerate}[label=(\roman*)]
  \item \textbf{Key difference:} BQP permits a constant bounded error; EQP requires exact correctness.
    BQP is believed to be a proper superset of EQP.
  \item \textbf{Example in BQP not known in P:} \textbf{Integer factoring}
    (Shor's algorithm, polynomial time on a quantum computer; no polynomial classical algorithm is known).
  \item \textbf{EQP is less robust} because: (a) EQP depends on the chosen universal gate set ---
    changing the gate set can alter which functions fall in EQP, making it gate-set-dependent;
    (b) exact-zero-error computations are extremely sensitive to amplitude precision and circuit
    symmetries, making EQP an artificial class for physical hardware.
    BQP, by contrast, is gate-set-independent and matches the notion of ``efficiently quantum-computable''
    under any standard error model.
\end{enumerate}

\subsection*{(c) QPSPACE $=$ PSPACE [5 pts]}

\begin{enumerate}[label=(\roman*)]
  \item \textbf{PSPACE:} All decision problems solvable by a deterministic classical computer
    using polynomial memory (with no time restriction).\\
    \textbf{QPSPACE:} All decision problems solvable by a quantum computer
    using polynomial memory (with no time restriction).

  \item \textbf{Implication:} QPSPACE$=$PSPACE means quantum computers offer \textbf{no memory advantage}
    over classical computers. Any problem a quantum machine can solve in polynomial space,
    a classical machine can also solve in polynomial space, and vice versa.
    Contrast with \emph{time}: Shor's algorithm (factoring) and Grover's search provide exponential and
    quadratic time speedups respectively, but these are purely time-based advantages.

  \item \textbf{Intuition:} A quantum state on $n$ qubits has $2^n$ amplitudes, which naively requires
    exponential space. However, we can compute output amplitudes \emph{one at a time}:
    \[
    \alpha_{x_\text{out}}=\sum_{\text{paths}}\prod_{\text{gates}}\text{(matrix entry)}.
    \]
    Enumerating each path and accumulating the weighted sum requires only $O(n)$ space for
    the current path and running total, never storing the full $2^n$-dimensional state vector.
    Total space: $O(\text{poly}(n))$ for the circuit description plus $O(n)$ for path tracking.
    This shows QPSPACE$\subseteq$PSPACE; the reverse inclusion is trivial.
\end{enumerate}

\subsection*{(d) NISQ era and formal methods [5 pts]}

\begin{enumerate}[label=(\roman*)]
  \item \textbf{NISQ} = \textbf{Noisy Intermediate-Scale Quantum.}

    Two fundamental limitations:
    \begin{itemize}
      \item \emph{Noise and decoherence:} Gate errors ($\sim0.1$--$1\%$ per gate) and
        qubit decoherence times ($\sim$ microseconds) restrict usable circuit depth.
        Full fault-tolerance requires $\sim10^3$ physical qubits per logical qubit,
        far beyond current hardware.
      \item \emph{Limited scale and connectivity:} Current devices have 50--1000 physical qubits
        with sparse qubit-to-qubit connections, introducing routing overhead that increases
        effective circuit depth and error accumulation.
    \end{itemize}

  \item \textbf{WMC-based simulation:}
    Each gate is encoded as a Boolean constraint on per-wire-per-timestep variables.
    Reversible gates (NOT, CNOT, Toffoli) contribute equality constraints; the Hadamard
    gate introduces an auxiliary variable $h$ with weights $W(h)=-\tfrac{1}{\sqrt{2}}$,
    $W(\neg h)=+\tfrac{1}{\sqrt{2}}$ to encode the negative amplitude entry.
    The amplitude of output string $x_\text{out}$ is then:
    \[
    \alpha_{x_\text{out}}=\mathrm{WMC}_W\bigl(F_\text{circuit}\land(\text{output}=x_\text{out})\bigr).
    \]
    A weighted model counter exploits the product structure of the formula to compute exact
    amplitudes without storing the full $2^n$-dimensional state vector, enabling practical
    classical verification of small quantum circuits.

  \item \textbf{Feynman's connection:}
    Feynman proposed quantum computers (1982) precisely because classical computers are
    \emph{exponentially inefficient} at simulating quantum systems (computing ground-state
    energies, correlation functions, etc.).
    The hardness of verifying a quantum circuit classically --- checking that an optimized
    circuit computes the same unitary as the original, or computing a specific output amplitude
    --- inherits the same BQP and QMA hardness as the underlying physical problems.
    Thus the formal-methods tools needed to verify quantum circuits face exactly the computational
    barrier that motivated quantum computing in the first place. This explains both why:
    \begin{itemize}
      \item Classical verification tools (\#SAT, decision diagrams, ZX-calculus) cannot scale
        efficiently to large circuits without solving BQP $\stackrel{?}{=}$ P; and
      \item Studying circuit hardness via complexity theory remains fundamental even before
        fault-tolerant quantum computers are built.
    \end{itemize}
\end{enumerate}

\end{document}
